\documentclass[../readings.tex]{subfiles}

\begin{document}

\subsection{Lanczos Algorithm}
\label{sub:lanczos-theory}

Iterative method for finding extreme eigenpairs of large, symmetric matrices without explicit $O(n^2)$ formation.

\subsubsection{Krylov Subspaces}

\addDef{Krylov Subspace}{%
$\mathcal{K}_k(\bA, \bv) = \text{span}\{\bv,\, \bA\bv,\, \bA^2\bv,\, \dots,\, \bA^{k-1}\bv\}$.
\textbf{Intuition:}\enspace $\bA^j\bv$ amplifies components along dominant eigenvectors, concentrating information about the spectral edges.
}
\label{def:krylov-subspace}

\subsubsection{The Lanczos Recurrence}

\addThm{Lanczos Relation}{%
The process builds an orthonormal basis $V_k = [\bv_1, \dots, \bv_k]$ for $\mathcal{K}_k$ satisfying:
\[
    \bA V_k = V_k T_k + \beta_k \bv_{k+1} \be_k^T.
\]
\begin{itemize}
    \item \textbf{$T_k$:}\enspace Symmetric \textbf{tridiagonal} matrix.
    \item \textbf{$\alpha_j = \bv_j^T \bA \bv_j$:}\enspace Rayleigh quotients (diagonal).
    \item \textbf{$\beta_j = \|\text{residual}\|$:}\enspace Orthogonalization coefficients (off-diagonal).
\end{itemize}
}
\label{thm:lanczos-relation}

\addNote{%
\textbf{Performance Advantage:}\enspace
Unlike Gram-Schmidt, which orthogonalizes against all previous vectors, Lanczos only requires the \textbf{two previous vectors} (three-term recurrence). This makes it $O(k \cdot \text{nnz})$ for sparse matrices.
}

\subsubsection{Ritz Values and Vectors}

Approximations to the true eigenpairs from the Krylov subspace.

\addDef{Ritz Pair}{%
\begin{itemize}
    \item \textbf{Ritz Value ($\theta_i$):}\enspace Eigenvalue of $T_k$.
    \item \textbf{Ritz Vector ($\tilde{\bu}_i$):}\enspace $V_k \bs_i$, where $\bs_i$ is an eigenvector of $T_k$.
\end{itemize}
}
\label{def:ritz-pair}

\addThm{Residual Estimate}{%
The residual $\|\bA\tilde{\bu}_i - \theta_i \tilde{\bu}_i\|_2 = \beta_k |s_i(k)|$.
\textbf{Significance:}\enspace Convergence of a Ritz pair can be monitored for $O(k)$ cost without forming the full $n$-vector $\tilde{\bu}_i$.
}
\label{thm:ritz-residual}

\addNote{%
\textbf{Ghost Eigenvalues:}\enspace
In floating-point, the three-term recurrence loses orthogonality, leading to ``ghost'' (spurious) copies of eigenvalues. \textbf{Practical Rule:}\enspace Use full or selective reorthogonalization to maintain basis integrity.
}

\addExcr{%
\begin{enumerate}
    \item Carry out 2 steps of Lanczos for $\bA = \text{diag}(5, 3, 1)$ starting with $\bv_1 = \frac{1}{\sqrt{3}}\mathbf{1}$.
    \item Show that $T_k = V_k^T \bA V_k$ (Galerkin Projection).
    \item Implement Lanczos on a $50 \times 50$ Hilbert matrix. Count spurious eigenvalues without reorthogonalization.
\end{enumerate}
}

\end{document}
